% Supplementary Materials for
% "A Scoping Review of Visual Perception Payloads on Industrial Quadruped
% Robots: Sensing Modalities, Onboard Compute, and Reporting Completeness".
% Compile from this folder: pdflatex supplementary.tex (twice).
\documentclass[10pt]{article}
\usepackage[a4paper,margin=2cm]{geometry}
\usepackage[T1]{fontenc}
\usepackage{lmodern}
\usepackage{booktabs,tabularx,longtable,array}
\usepackage[hidelinks]{hyperref}
\usepackage{url}
\newcolumntype{C}{>{\centering\arraybackslash}X}
\newlength{\fulllength}\setlength{\fulllength}{\textwidth}
% numbers.tex -- GENERATED by analysis/scripts/40_emit_numbers.py. Do not edit.
% generated 2026-10-28T19:22:26 from: prisma_counts.csv, coding_time.csv, results_numbers.csv, pipeline_numbers.csv, review_time_numbers.csv
% A \numTODO value is one the analysis has not produced yet.
\newcommand{\Nidentified}{2{,}212}
\newcommand{\NidIEEE}{729}
\newcommand{\NidScopus}{1{,}137}
\newcommand{\NidWoS}{346}
\newcommand{\Ndupwithin}{575}
\newcommand{\Nbaselinematch}{88}
\newcommand{\Nscreened}{1{,}549}
\newcommand{\NexclTA}{1{,}464}
\newcommand{\NexclTAone}{1{,}459}
\newcommand{\NexclTAtwo}{14}
\newcommand{\NexclTAmodel}{1{,}473}
\newcommand{\NreinstatedTA}{9}
\newcommand{\Nsought}{85}
\newcommand{\Nnotretrieved}{18}
\newcommand{\Nassessed}{67}
\newcommand{\NexclFT}{47}
\newcommand{\Nincludednew}{20}
\newcommand{\Nprevious}{88}
\newcommand{\Nincluded}{108}
\newcommand{\NcodingStudies}{23}
\newcommand{\CodingTotalH}{14.0}
\newcommand{\CodingMedianMin}{35.4}
\newcommand{\CodingIQRlo}{23.4}
\newcommand{\CodingIQRhi}{44.3}
\newcommand{\CodingMinMin}{18.7}
\newcommand{\CodingMaxMin}{70.6}
\newcommand{\TimeCodebookMin}{43}
\newcommand{\TimeAdjMin}{48}
\newcommand{\TimeReadjMin}{73}
\newcommand{\TimeBaselineReviewMin}{16}
\newcommand{\TimeTitleScanMin}{135}
\newcommand{\ManualProjectedH}{63}
\newcommand{\PipelineHumanH}{2.7}
\newcommand{\TimeRatio}{23}
\newcommand{\TimeRatioWithVal}{3.8}
\newcommand{\NstudiesAll}{108}
\newcommand{\NfailsCurrent}{16}
\newcommand{\NroundOne}{88}
\newcommand{\NroundTwo}{20}
\newcommand{\YearFirst}{2019}
\newcommand{\YearLast}{2026}
\newcommand{\YearRecent}{74}
\newcommand{\YearRecentShare}{69}
\newcommand{\YearTwentyFour}{27}
\newcommand{\YearTwentyFive}{31}
\newcommand{\SecGeneral}{20}
\newcommand{\SecConstruction}{27}
\newcommand{\SecManufacturing}{7}
\newcommand{\SecOil}{13}
\newcommand{\SecNuclear}{11}
\newcommand{\SecEnergy}{14}
\newcommand{\SecOther}{5}
\newcommand{\SecMining}{14}
\newcommand{\CompXeightsixMiniPc}{13}
\newcommand{\CompEmbeddedGpu}{20}
\newcommand{\CompSbcMicrocontroller}{8}
\newcommand{\CompLaptopWorkstation}{10}
\newcommand{\CompVendorIntegrated}{5}
\newcommand{\CompNone}{57}
\newcommand{\SecMulti}{3}
\newcommand{\EnvReal}{59}
\newcommand{\EnvBoth}{38}
\newcommand{\EnvSim}{5}
\newcommand{\EnvUnclear}{5}
\newcommand{\VenueConferencePct}{52}
\newcommand{\CiteInstrument}{\cite{inc:paper001,inc:paper016,inc:paper027,inc:paper046,inc:paper048,inc:paper055,inc:paper062,inc:paper071,inc:r2-1437}}
\newcommand{\CiteTracking}{\cite{inc:paper020,inc:paper028,inc:paper030,inc:paper052,inc:paper074,inc:r2-0553,inc:r2-0934}}
\newcommand{\CiteEvent}{\cite{inc:paper059}}
\newcommand{\CiteArtefact}{\cite{inc:paper008,inc:paper013,inc:paper017,inc:paper032,inc:paper057,inc:paper059,inc:paper073,inc:paper074,inc:r2-0740,inc:r2-0934,inc:r2-1156}}
\newcommand{\CompMultiGpuPct}{43}
\newcommand{\CompSingleGpuPct}{36}
\newcommand{\CompNamed}{50}
\newcommand{\CompGpuAmongNamed}{20}
\newcommand{\CompMultiNamed}{28}
\newcommand{\CompMultiGpu}{12}
\newcommand{\CompSingleNamed}{22}
\newcommand{\CompSingleGpu}{8}
\newcommand{\HolisticN}{87}
\newcommand{\HolisticYes}{20}
\newcommand{\HolisticYesRCI}{6}
\newcommand{\HolisticNoRCI}{5}
\newcommand{\HolisticYesArtefact}{5}
\newcommand{\HolisticYesNoSensorModel}{8}
\newcommand{\HolisticYesNoComputeModel}{7}
\newcommand{\VenueConference}{56}
\newcommand{\SConstrLidar}{63}
\newcommand{\SConstrMono}{56}
\newcommand{\SConstrRgbd}{30}
\newcommand{\SConstrThermal}{11}
\newcommand{\SConstrStereo}{4}
\newcommand{\SConstrEvent}{4}
\newcommand{\SGenLidar}{45}
\newcommand{\SGenMono}{40}
\newcommand{\SGenRgbd}{40}
\newcommand{\SGenThermal}{25}
\newcommand{\SGenStereo}{5}
\newcommand{\SGenEvent}{0}
\newcommand{\SEnergyLidar}{50}
\newcommand{\SEnergyMono}{71}
\newcommand{\SEnergyRgbd}{36}
\newcommand{\SEnergyThermal}{57}
\newcommand{\SEnergyStereo}{7}
\newcommand{\SEnergyEvent}{0}
\newcommand{\SMiningLidar}{100}
\newcommand{\SMiningMono}{21}
\newcommand{\SMiningRgbd}{21}
\newcommand{\SMiningThermal}{36}
\newcommand{\SMiningStereo}{14}
\newcommand{\SMiningEvent}{0}
\newcommand{\SOilLidar}{46}
\newcommand{\SOilMono}{23}
\newcommand{\SOilRgbd}{54}
\newcommand{\SOilThermal}{15}
\newcommand{\SOilStereo}{0}
\newcommand{\SOilEvent}{0}
\newcommand{\SNuclearLidar}{36}
\newcommand{\SNuclearMono}{27}
\newcommand{\SNuclearRgbd}{18}
\newcommand{\SNuclearThermal}{0}
\newcommand{\SNuclearStereo}{0}
\newcommand{\SNuclearEvent}{0}
\newcommand{\SManufLidar}{43}
\newcommand{\SManufMono}{43}
\newcommand{\SManufRgbd}{71}
\newcommand{\SManufThermal}{0}
\newcommand{\SManufStereo}{43}
\newcommand{\SManufEvent}{0}
\newcommand{\SOtherLidar}{40}
\newcommand{\SOtherMono}{60}
\newcommand{\SOtherRgbd}{80}
\newcommand{\SOtherThermal}{0}
\newcommand{\SOtherStereo}{20}
\newcommand{\SOtherEvent}{0}
\newcommand{\ModLiDAR}{60}
\newcommand{\ModMonocular}{47}
\newcommand{\ModRGBD}{40}
\newcommand{\ModThermal}{23}
\newcommand{\ModStereo}{9}
\newcommand{\ModEvent}{1}
\newcommand{\NreadableStudies}{107}
\newcommand{\ModNone}{11}
\newcommand{\OtherCutoff}{7}
\newcommand{\TopComboName}{a LiDAR}
\newcommand{\TopComboN}{22}
\newcommand{\ModAny}{96}
\newcommand{\ModMulti}{53}
\newcommand{\ModMultiPctOfAny}{55}
\newcommand{\LidarWithCamera}{38}
\newcommand{\LidarAlone}{22}
\newcommand{\ModOne}{43}
\newcommand{\ModTwo}{29}
\newcommand{\ModThreePlus}{24}
\newcommand{\ModCombos}{22}
\newcommand{\FigFourDropped}{Event camera (1)}
\newcommand{\XLidarSlam}{72}
\newcommand{\XLidarNav}{83}
\newcommand{\XLidarObst}{55}
\newcommand{\XLidarDet}{48}
\newcommand{\XLidarInstr}{8}
\newcommand{\XLidarTrack}{3}
\newcommand{\XMonoSlam}{55}
\newcommand{\XMonoNav}{77}
\newcommand{\XMonoObst}{53}
\newcommand{\XMonoDet}{55}
\newcommand{\XMonoInstr}{13}
\newcommand{\XMonoTrack}{9}
\newcommand{\XRgbdSlam}{65}
\newcommand{\XRgbdNav}{72}
\newcommand{\XRgbdObst}{60}
\newcommand{\XRgbdDet}{62}
\newcommand{\XRgbdInstr}{5}
\newcommand{\XRgbdTrack}{10}
\newcommand{\XThermalSlam}{52}
\newcommand{\XThermalNav}{87}
\newcommand{\XThermalObst}{52}
\newcommand{\XThermalDet}{65}
\newcommand{\XThermalInstr}{4}
\newcommand{\XThermalTrack}{4}
\newcommand{\XStereoSlam}{78}
\newcommand{\XStereoNav}{89}
\newcommand{\XStereoObst}{67}
\newcommand{\XStereoDet}{44}
\newcommand{\XStereoInstr}{0}
\newcommand{\XStereoTrack}{11}
\newcommand{\XEventSlam}{0}
\newcommand{\XEventNav}{0}
\newcommand{\XEventObst}{100}
\newcommand{\XEventDet}{100}
\newcommand{\XEventInstr}{0}
\newcommand{\XEventTrack}{0}
\newcommand{\XNoneSlam}{0}
\newcommand{\XNoneNav}{55}
\newcommand{\XNoneObst}{45}
\newcommand{\XNoneDet}{0}
\newcommand{\XNoneInstr}{0}
\newcommand{\XNoneTrack}{9}
\newcommand{\CapSLAM}{58}
\newcommand{\CapAutonomous}{76}
\newcommand{\CapObstacle}{53}
\newcommand{\CapObject}{48}
\newcommand{\CapInstrument}{9}
\newcommand{\CapTarget}{7}
\newcommand{\CapRealWorld}{82}
\newcommand{\CapGeometricAny}{88}
\newcommand{\CapSemanticAny}{52}
\newcommand{\CapGeometricOnly}{44}
\newcommand{\CapSemanticOnly}{8}
\newcommand{\CapNeither}{11}
\newcommand{\SensSpecUnspecified}{11}
\newcommand{\SensSpecModality}{43}
\newcommand{\SensSpecModel}{53}
\newcommand{\CompSpecUnspecified}{57}
\newcommand{\CompSpecModality}{6}
\newcommand{\CompSpecModel}{44}
\newcommand{\RCISectorLowName}{energy and power}
\newcommand{\RCISectorLowMedian}{3.5}
\newcommand{\RCISectorHighMedian}{6}
\newcommand{\NrciStudies}{107}
\newcommand{\RCImedian}{5}
\newcommand{\RCIiqrLo}{4}
\newcommand{\RCIiqrHi}{6}
\newcommand{\RCImin}{0}
\newcommand{\RCImax}{8}
\newcommand{\RCIfull}{0}
\newcommand{\RCIatMostThree}{18}
\newcommand{\RCIthreeStrict}{31}
\newcommand{\RCIitemOne}{82}
\newcommand{\RCIitemTwo}{53}
\newcommand{\RCIitemThree}{44}
\newcommand{\RCIitemFour}{15}
\newcommand{\RCIitemFive}{81}
\newcommand{\RCIitemSix}{96}
\newcommand{\RCIitemSeven}{78}
\newcommand{\RCIitemEight}{82}
\newcommand{\RCIitemNine}{11}
\newcommand{\NsensStudies}{91}
\newcommand{\SensMaxDiffPP}{8}
\newcommand{\SensMaxDiffItem}{success metric (RCI item 8)}
\newcommand{\SensRCImedian}{5}
\newcommand{\PipeCodes}{368}
\newcommand{\PipeWrong}{49}
\newcommand{\PipeWrongPct}{13}
\newcommand{\PipeAgreed}{334}
\newcommand{\PipeAgreedWrong}{34}
\newcommand{\PipeAgreedWrongPct}{10}
\newcommand{\PipeAgreedFP}{25}
\newcommand{\PipeAdj}{34}
\newcommand{\PipeAdjMatch}{19}
\newcommand{\PipeAdjWrong}{15}
\newcommand{\PipeAdjWrongToYes}{14}
\newcommand{\PipeYes}{183}
\newcommand{\HumanYes}{154}
\newcommand{\PipeYesInflationPct}{19}
\newcommand{\ReadjN}{67}
\newcommand{\ReadjToNo}{21}
\newcommand{\ReadjToNoPct}{31}
\newcommand{\InterAgreeMin}{84}
\newcommand{\InterAgreeMax}{99}
\newcommand{\InterKappaMin}{0.45}
\newcommand{\InterKappaMax}{0.97}
\newcommand{\InterKappaLowItems}{rci4 middleware version, rci8 success metric, cap real world validation}
\newcommand{\HumanAgreeMin}{65}
\newcommand{\HumanAgreeMax}{100}
\newcommand{\HumanAgreeLowItem}{real world validation}
\newcommand{\HumanAgreeLowValue}{65}
\newcommand{\ClassNoneAgreeSensors}{83}
\newcommand{\ClassNoneAgreeCompute}{91}
\newcommand{\ClassSensAgreeMin}{65}
\newcommand{\ClassSensAgreeMax}{100}
\newcommand{\ClassCompAgreeMin}{96}
\newcommand{\ClassCompAgreeMax}{100}
\newcommand{\ClassStudies}{23}
\newcommand{\ClassMonoAgree}{65}
\newcommand{\ClassMonoPipe}{10}
\newcommand{\ClassMonoHuman}{4}
\newcommand{\SpecWkSens}{0.63}
\newcommand{\SpecWkComp}{0.73}
\newcommand{\RecallN}{88}
\newcommand{\RecallKept}{82}
\newcommand{\RecallPct}{93}
\newcommand{\RecallLo}{86}
\newcommand{\RecallHi}{97}
\newcommand{\RecallEligN}{72}
\newcommand{\RecallEligKept}{70}
\newcommand{\RecallEligPct}{97}
\newcommand{\RecallEligLo}{90}
\newcommand{\RecallEligHi}{99}
\newcommand{\RecallLostFlagged}{4}
\newcommand{\RecallLost}{6}

\renewcommand{\thetable}{S\arabic{table}}
% The shared reference file carries author-year labels for the MDPI class;
% here the references are numbered, so the optional label is dropped.
\let\origbibitem\bibitem
\renewcommand{\bibitem}[2][]{\origbibitem{#2}}
\setlength{\parindent}{0pt}\setlength{\parskip}{4pt}

\title{Supplementary Materials\\[2pt]\large A Scoping Review of Visual
Perception Payloads on Industrial Quadruped Robots}
\author{Panna Zsoldos and Péter Galambos}
\date{}

\begin{document}
\maketitle

\noindent Contents: Table~S1, PRISMA-ScR checklist. Table~S2, included
studies with their charted characteristics and references. Table~S3,
deviations from the registered protocol. Section~S1, further technical
details of the charting pipeline. Table~S4, extraction accuracy in detail.
Table~S5, class-level validation of the controlled vocabulary. Table~S6,
sensitivity analysis. Files S1--S4 (prompts, codebooks, adjudication logs,
reports not retrieved) are part of the data release referenced in the Data
Availability Statement.

%-----------------------------------------------------------------
\section*{Table S1. PRISMA-ScR checklist}

\begin{longtable}{p{0.3cm}p{4.2cm}p{11.2cm}}
\toprule
\# & Item & Reported in\\
\midrule\endhead
\bottomrule\endfoot
1 & Title & Title: identified as a scoping review.\\
2 & Structured summary & Abstract.\\
3 & Rationale & Section 1, paragraphs 1--4 (payload as the deciding factor; gap in existing reviews; AXIS feasibility study).\\
4 & Objectives & Section 1: aim, PCC elements and RQ1--RQ4.\\
5 & Protocol and registration & Section 3.1: OSF registration (https://osf.io/r8zs6), partly retrospective; thirteen deviations, Table~S3.\\
6 & Eligibility criteria & Section 3.4 and Table 1 (PCC); language and year limits; excluded publication types; application of the context criterion.\\
7 & Information sources & Section 3.3: Scopus, IEEE Xplore, Web of Science Core Collection; search date 3 September 2026; earlier review as previous-studies source.\\
8 & Search & Appendix A, Table A1: full strings as executed for all three databases, with filters and hits.\\
9 & Selection of sources of evidence & Sections 3.5 and 3.9: two-stage LLM-assisted screening, human title scan of all exclusions, human full-text decisions with advisory model passes, recall check on the earlier studies, audit of the earlier studies.\\
10 & Data charting process & Sections 3.2, 3.6 and 3.9: two independent LLM passes with verified quotations, human adjudication, re-adjudication with context, blind human validation subsample; Section~S1.\\
11 & Data items & Section 3.6 (sensors, compute, capabilities, real-world validation, sector, environment), Section 3.7 (specificity) and Section 3.8 with Table 3 (RCI items).\\
12 & Critical appraisal of individual sources & Not performed, as is usual for scoping reviews; the RCI charts reporting completeness, not study quality (Sections 3.8 and 6).\\
13 & Synthesis of results & Section 3.10: descriptive frequencies and cross-tabulations; sensitivity analysis (Table~S6).\\
14 & Selection of sources of evidence & Section 4.1 and Figure 1 (flow diagram with numbers and exclusion reasons).\\
15 & Characteristics of sources of evidence & Section 4.1 and Table~S2 (characteristics and citation of every included study).\\
16 & Critical appraisal within sources & Not applicable (item 12).\\
17 & Results of individual sources of evidence & Table~S2; full charting data in the data release.\\
18 & Synthesis of results & Sections 4.2--4.6, Tables 2, 4 and 5, Figures 2--6.\\
19 & Summary of evidence & Section 5.1.\\
20 & Limitations & Section 6.\\
21 & Conclusions & Section 7.\\
22 & Funding & Funding statement.\\
\end{longtable}

%-----------------------------------------------------------------
\section*{Table S2. Included studies}

Sector: Gen., industrial setting without a named sector; Constr.,
construction; Oil/gas/chem., oil, gas and chemical plants; Manuf.,
manufacturing. Modalities: L, LiDAR; C, monocular camera; D, RGB-D or depth
camera; T, thermal or infrared camera; S, stereo camera; E, event camera.
Compute: GPU, embedded GPU module; x86, small-form-factor x86 computer; CPU,
laptop- or desktop-class processor; SBC, single-board computer or
microcontroller; Vendor, vendor-integrated computer. Capabilities: M, SLAM or
mapping; N, autonomous navigation; O, obstacle avoidance; D, object
detection; I, instrument reading; T, target tracking. RCI, Reporting
Completeness Index (0--9). $\dagger$, an earlier study that would not meet
the eligibility criteria as applied in round~2. A dash means not reported;
the scanned study has no codes. Reference numbers refer to the list at the
end of this document.

{\footnotesize
\begin{longtable}{lllllll}
\toprule
\textbf{Study} & \textbf{Year} & \textbf{Sector} & \textbf{Modalities} & \textbf{Compute} & \textbf{Capabilities} & \textbf{RCI}\\
\midrule
\endhead
\bottomrule
\endfoot
Wisth et al.~\cite{inc:paper014}$\dagger$ & 2019 & Gen. & LCDT & -- & NO & 6\\
Ramezani et al.~\cite{inc:paper019} & 2020 & Oil/gas/chem. & L & -- & MN & 6\\
Wang et al.~\cite{inc:paper040} & 2020 & Gen. & L & -- & NOD & 6\\
Miller et al.~\cite{inc:paper005} & 2020 & Mining & LDTS & GPU & MNOD & 7\\
Xin et al.~\cite{inc:paper007} & 2020 & Gen. & CT & CPU & N & 5\\
Brandao et al.~\cite{inc:paper073} & 2020 & Oil/gas/chem. & L & -- & MNO & 5\\
Buchanan et al.~\cite{inc:r2-0477} & 2020 & Mining & L & -- & MN & 5\\
Bouman et al.~\cite{inc:r2-0192} & 2020 & Energy & LCT & -- & MNOD & 5\\
Kim et al.~\cite{inc:paper065} & 2021 & Constr. & L & GPU & MD & 7\\
Jud et al.~\cite{inc:paper087} & 2021 & Constr. & LC & CPU & N & 7\\
Dong et al.~\cite{inc:paper075}$\dagger$ & 2021 & Energy & LCDT & -- & MNOD & 1\\
Wu et al.~\cite{inc:paper036} & 2021 & Energy & DT & -- & MNOD & 1\\
Mattamala et al.~\cite{inc:paper038} & 2021 & Gen. & CD & CPU & MN & 7\\
Stachowiak et al.~\cite{inc:paper078} & 2021 & Mining & LT & -- & ND & 4\\
Skoczylas et al.~\cite{inc:paper085} & 2021 & Mining & LCT & -- & N & 4\\
Gehring et al.~\cite{inc:r2-1437} & 2021 & Energy & LCDT & -- & MNODI & 6\\
Halder et al.~\cite{inc:paper069} & 2022 & Constr. & C & SBC & O & 5\\
Wetzel et al.~\cite{inc:paper070} & 2022 & Constr. & LC & Vendor & NO & 6\\
Ding et al.~\cite{inc:paper018}$\dagger$ & 2022 & Energy & -- & -- & N & 0\\
Kulkarni et al.~\cite{inc:paper013} & 2022 & Mining & L & CPU & MNO & 7\\
Mattamala et al.~\cite{inc:paper025} & 2022 & Mining & LCS & CPU & MNO & 7\\
Tang et al.~\cite{inc:paper067}$\dagger$ & 2022 & Mining & LD & -- & MN & 4\\
Lysakowski et al.~\cite{inc:r2-0387} & 2022 & Mining & L & -- & MN & 5\\
Jang et al.~\cite{inc:r2-1156} & 2022 & Nuclear & D & -- & MD & 5\\
He et al.~\cite{inc:paper026}$\dagger$ & 2023 & Manuf. & LDS & GPU, SBC & MN & 3\\
Molnar et al.~\cite{inc:paper084}$\dagger$ & 2023 & Oil/gas/chem. & -- & -- & NO & 4\\
Halder et al.~\cite{inc:paper066} & 2023 & Constr. & LC & -- & MND & 6\\
Fiorucci et al.~\cite{inc:paper080} & 2023 & Energy & LCDT & CPU & MND & 7\\
Li et al.~\cite{inc:paper077} & 2023 & Energy & C & -- & -- & 5\\
Zhang et al.~\cite{inc:paper027} & 2023 & Constr. & C & -- & MNDI & 3\\
Ly et al.~\cite{inc:paper044} & 2023 & Oil/gas/chem. & D & -- & MNO & 6\\
Crespo et al.~\cite{inc:paper064} & 2023 & Mining & LT & -- & NO & 4\\
Chu et al.~\cite{inc:paper072} & 2023 & Constr. & CS & GPU & MNOD & 7\\
Brogaard et al.~\cite{inc:r2-0495} & 2023 & Oil/gas/chem. & CD & GPU, x86 & D & 7\\
Lee et al.~\cite{inc:paper021} & 2024 & Oil/gas/chem. & D & GPU, x86 & NOD & 4\\
Lee et al.~\cite{inc:paper068} & 2024 & Oil/gas/chem. & -- & -- & NO & 4\\
Dassori et al.~\cite{inc:paper006}$\dagger$ & 2024 & Constr. & C & -- & NO & 4\\
Gan et al.~\cite{inc:paper086} & 2024 & Constr. & L & GPU & MN & 3\\
Baru A~\cite{inc:paper079} & 2024 & Constr. & C & -- & NOD & 4\\
Wu et al.~\cite{inc:paper063} & 2024 & Oil/gas/chem. & L & -- & MNO & 5\\
Tsenis et al.~\cite{inc:paper081} & 2024 & Energy & LCT & -- & N & 0\\
Aziz et al.~\cite{inc:paper062} & 2024 & Constr. & C & CPU, Vendor & MNDI & 5\\
Ginting et al.~\cite{inc:paper071}$\dagger$ & 2024 & Constr. & LC & -- & NODI & 4\\
Hu et al.~\cite{inc:paper083} & 2024 & Constr. & LDT & GPU & MD & 4\\
Rowell et al.~\cite{inc:paper041} & 2024 & Gen. & L & -- & MND & 6\\
Chang et al.~\cite{inc:paper012} & 2024 & Gen. & T & -- & -- & 1\\
Zhou et al.~\cite{inc:paper015} & 2024 & Gen. & LC & -- & MN & 4\\
Reese et al.~\cite{inc:paper037} & 2024 & Gen. & D & -- & D & 3\\
Zagorski et al.~\cite{inc:paper082} & 2024 & Gen. & L & -- & -- & 3\\
Hou et al.~\cite{inc:paper039} & 2024 & Constr., Energy & L & SBC & -- & 2\\
Turrisi et al.~\cite{inc:paper074} & 2024 & Constr. & -- & -- & NT & 4\\
Wei et al.~\cite{inc:paper033}$\dagger$ & 2024 & Manuf. & LDS & -- & MNO & 4\\
Sun et al.~\cite{inc:paper076}$\dagger$ & 2024 & Mining & LD & -- & MNOD & 5\\
Chen et al.~\cite{inc:paper048} & 2024 & Mining & L & -- & DI & 7\\
Balan et al.~\cite{inc:paper035} & 2024 & Nuclear & LC & x86 & MNOD & 6\\
Xu et al.~\cite{inc:paper049} & 2024 & Nuclear & C & -- & NO & 5\\
Cheng et al.~\cite{inc:paper023} & 2024 & Gen. & L & x86 & MN & 7\\
Li et al.~\cite{inc:r2-0771} & 2024 & Nuclear & -- & -- & NO & 5\\
Baird et al.~\cite{inc:r2-1052} & 2024 & Gen. & T & Vendor & MNOD & 6\\
Huang et al.~\cite{inc:r2-1239} & 2024 & Energy & C & CPU & MO & 3\\
Moon et al.~\cite{inc:r2-1169} & 2024 & Gen. & CDT & -- & N & 2\\
Damigos et al.~\cite{inc:paper042} & 2025 & Gen. & DS & x86 & -- & 5\\
Gan et al.~\cite{inc:paper047} & 2025 & Constr. & L & -- & MNO & 5\\
Zeng et al.~\cite{inc:paper003} & 2025 & Constr. & CD & GPU & D & 6\\
Scheidemann et al.~\cite{inc:paper030}$\dagger$ & 2025 & Gen. & LD & GPU & MNODT & 6\\
Chung et al.~\cite{inc:paper053} & 2025 & Constr. & LD & x86 & MNOD & 7\\
Stührenberg et al.~\cite{inc:paper057} & 2025 & Constr. & LD & x86 & M & 8\\
Kim et al.~\cite{inc:paper061} & 2025 & Constr. & L & x86 & D & 6\\
Xu et al.~\cite{inc:paper011} & 2025 & Nuclear & -- & -- & NO & 4\\
Kim et al.~\cite{inc:paper058} & 2025 & Other & CD & x86 & MO & 6\\
Xing et al.~\cite{inc:paper028} & 2025 & Energy & CDTS & GPU & MNODT & 5\\
Mu et al.~\cite{inc:paper043} & 2025 & Gen. & -- & -- & -- & --\\
Hu et al.~\cite{inc:paper010} & 2025 & Gen. & -- & -- & -- & 4\\
Yang et al.~\cite{inc:paper016} & 2025 & Energy & LC & -- & MNODI & 3\\
Zhu et al.~\cite{inc:paper020}$\dagger$ & 2025 & Manuf. & C & GPU & MNT & 7\\
Mendoza et al.~\cite{inc:paper024}$\dagger$ & 2025 & Gen. & C & -- & -- & 3\\
Khatiri et al.~\cite{inc:paper031} & 2025 & Gen. & LCD & -- & NO & 5\\
Rizvi et al.~\cite{inc:paper029} & 2025 & Manuf. & DS & -- & MND & 5\\
DiBattista et al.~\cite{inc:paper051} & 2025 & Other & LC & GPU & MNO & 5\\
Fan et al.~\cite{inc:paper004} & 2025 & Mining & L & -- & MNO & 2\\
Gao et al.~\cite{inc:paper017}$\dagger$ & 2025 & Mining & LCT & -- & MND & 5\\
Li et al.~\cite{inc:paper009} & 2025 & Nuclear & -- & -- & -- & 3\\
Staniaszek et al.~\cite{inc:paper045} & 2025 & Nuclear, Oil/gas/chem. & LC & x86 & MN & 6\\
Liu et al.~\cite{inc:paper050} & 2025 & Nuclear & -- & SBC & O & 4\\
Yao et al.~\cite{inc:paper001} & 2025 & Oil/gas/chem. & D & GPU & I & 5\\
Bin Mohamad Nor et al.~\cite{inc:paper002} & 2025 & Oil/gas/chem. & CDT & CPU & MNOD & 7\\
Guo et al.~\cite{inc:paper034} & 2025 & Oil/gas/chem. & LD & -- & NO & 6\\
Wu et al.~\cite{inc:paper054} & 2025 & Gen. & -- & -- & -- & 4\\
Zhang et al.~\cite{inc:paper060} & 2025 & Constr. & L & GPU & MD & 8\\
Shi et al.~\cite{inc:paper022}$\dagger$ & 2025 & Other & LCD & -- & MNOD & 1\\
Adams et al.~\cite{inc:r2-0704} & 2025 & Nuclear & LD & Vendor & MNO & 6\\
Chen et al.~\cite{inc:r2-0745} & 2025 & Constr. & LCD & GPU, x86, SBC & MNOD & 8\\
Tuomisto et al.~\cite{inc:paper056} & 2026 & Constr. & LCDT & -- & MNOD & 7\\
Lee et al.~\cite{inc:paper055} & 2026 & Energy & C & GPU & DI & 6\\
Liu et al.~\cite{inc:paper088} & 2026 & Energy & T & SBC & -- & 4\\
Ji et al.~\cite{inc:paper059} & 2026 & Constr. & CDE & x86 & OD & 8\\
Luo et al.~\cite{inc:paper032}$\dagger$ & 2026 & Other & DS & -- & NO & 7\\
Prajapati et al.~\cite{inc:paper008} & 2026 & Oil/gas/chem. & LDT & GPU & NOD & 8\\
Zhu et al.~\cite{inc:paper052} & 2026 & Manuf. & LCD & SBC & MODT & 5\\
Adams et al.~\cite{inc:paper046} & 2026 & Nuclear & L & CPU & MNOI & 7\\
Aydogmus et al.~\cite{inc:r2-0348} & 2026 & Gen. & CD & GPU & MND & 7\\
Ovur et al.~\cite{inc:r2-0553} & 2026 & Constr., Manuf., Other & D & Vendor & NDT & 8\\
Son et al.~\cite{inc:r2-0677} & 2026 & Nuclear & -- & GPU & -- & 6\\
Feng et al.~\cite{inc:r2-0740} & 2026 & Constr. & L & -- & MD & 7\\
Haage et al.~\cite{inc:r2-0776} & 2026 & Constr. & LCT & -- & ND & 5\\
Qin et al.~\cite{inc:r2-0934} & 2026 & Manuf. & C & -- & DT & 6\\
Bayer et al.~\cite{inc:r2-1004} & 2026 & Mining & L & x86 & MNO & 7\\
Chang et al.~\cite{inc:r2-1031} & 2026 & Gen. & L & SBC & MNOD & 6\\
\end{longtable}

}

%-----------------------------------------------------------------
\section*{Table S3. Deviations from the registered protocol}

Each deviation is logged with its full reasoning in the amendment log of
the OSF registration and in the data release.

{\small
\begin{longtable}{p{0.8cm}p{2.0cm}p{12.6cm}}
\toprule
& Date & Deviation\\
\midrule\endhead
\bottomrule\endfoot
A1 & 4 Sep 2026 & RCI item 9 restricted to artifacts released by the authors; publisher supplementary-material notices and third-party tools excluded. Pre-amendment codes retained.\\
A2 & 4 Sep 2026 & Documents with fewer than 1{,}000 characters of text per page are not coded by the models (one scanned study).\\
A3 & 26 Sep 2026 & Pass-B model replaced after the endpoint withdrew the registered model; pass B re-run over the whole corpus.\\
A4 & 26 Sep 2026 & Application of the context criterion at full text stated explicitly (laboratory validation of an industrial task qualifies).\\
A5 & 8 Oct 2026 & Validation subsample set at 24 studies (18 earlier, 6 round-2) before any agreement was computed; the first 18 drawn earlier studies in record order were kept.\\
A6 & 20 Oct 2026 & One drawn round-2 study excluded at full text before coding; subsample of 23, not replaced.\\
A7 & 20 Oct 2026 & Five human reference codes on RCI item 4 (ROS without a version) corrected to the rule after unblinding; blind codes retained.\\
A8 & 21 Oct 2026 & Human title scan of all records excluded by the models at title and abstract level; nine sent to full text, three included.\\
A9 & 19--21 Oct 2026 & Technical: decoding of a constant character offset in two PDFs; longer request timeout and token limit for pass B.\\
A10 & 23 Oct 2026 & Two rules fixed after codes had been seen: a product family counts as model level (strict version of item 3 reported); stairs are terrain, not obstacles.\\
A11 & 23 Oct 2026 & Disagreements first decided as positive re-adjudicated with context; blind human codes used as final codes for the validation studies.\\
A12 & 28 Oct 2026 & Eligibility audit of the 88 earlier studies; 16 that would not be included today kept and reported in a sensitivity analysis.\\
A13 & 28 Oct 2026 & One further human reference code corrected (RCI item 2 given for a radiation detector); sensor and compute strings taken from the human coding for the validation studies and from pass B where pass A listed none.\\
\end{longtable}
}

%-----------------------------------------------------------------
\section*{Section S1. Technical details of the charting pipeline}

\textbf{Reading.} Each document was converted to text with PyMuPDF and read
in overlapping windows of 30{,}000 characters with a 4{,}000-character
overlap. A positive code in any window counts for the study; where several
windows support it, the best-verified quotation is kept. A positive without
a verifiable quotation is not counted. Two documents whose embedded font
shifted every character by a constant offset were decoded before
extraction, after one of them had appeared to report nothing.

\textbf{Quotation matching.} Quotations are matched after lower-casing and
removing every character that is not a letter or digit, so that line breaks,
hyphenation and ligatures in the text layer do not cause false rejections.
Elided quotations are accepted only when every fragment is found and the
fragments cover most of the quoted text.

\textbf{Models and endpoint.} Both models were served by an institutional
Open WebUI/Ollama endpoint without a random seed. When checked on
26~October 2026, the endpoint reported the build identifiers b9371 for
gpt-oss:120b and b11332 for Qwen3.8-Flash-Next; it does not publish the
parameter count or quantization of the latter. Pass B used a 900~s request
timeout and a 16{,}000-token limit, because the model spends part of its
budget on reasoning. Every request, with its parameters and date, is logged
in the data release.

\textbf{Search execution.} In Scopus every term was kept quoted: unquoted
single words are expanded by the database and raised the count from
\NidScopus{} to 1{,}782 for otherwise identical logic. The classic Scopus
interface returns the larger count for the quoted string as well, so the
reported count refers to the current interface. In IEEE Xplore the
truncation wildcards inside quoted phrases were retained; without them the
count fell to 546.

\textbf{Screening recall.} The title and abstract procedure, re-run on the
\RecallN{} earlier studies, retained \RecallKept{} (\RecallPct\%). Of the
\RecallLost{} it would have excluded, \RecallLostFlagged{} are among the
earlier studies that would not be included today.

%-----------------------------------------------------------------
\section*{Table S4. Extraction accuracy in detail}

Agreement of each model pass with the blind human reference on the
\NcodingStudies{} validation studies, with Wilson 95\% intervals for the
agreement and asymptotic 95\% intervals for $\kappa$, clipped to $[-1,1]$.
With this sample size the $\kappa$ intervals are wide and approximate. Rows
marked as blind codes use the reviewer's codes before the corrections in
A7 and A13.

{\footnotesize
\setlength{\fulllength}{\textwidth}
\begin{tabularx}{\fulllength}{lcCCCCC}
\toprule
\textbf{Item} & \textbf{Pass} & \textbf{Agreement, \% (95\% CI)} & \textbf{Prev.\ human, \%} & \textbf{Prev.\ model, \%} & \textbf{$\kappa$ (95\% CI)} & \textbf{Model-only / human-only positives}\\
\midrule
1 Platform model & A & 91 (73--98) & 78 & 70 & 0.78 (0.48, 1.00) & 0 / 2\\
 & B & 96 (79--99) & 78 & 74 & 0.88 (0.65, 1.00) & 0 / 1\\
2 Sensor model & A & 96 (79--99) & 39 & 35 & 0.91 (0.73, 1.00) & 0 / 1\\
 & B & 96 (79--99) & 39 & 43 & 0.91 (0.74, 1.00) & 1 / 0\\
3 Compute model & A & 83 (63--93) & 35 & 35 & 0.62 (0.28, 0.96) & 2 / 2\\
 & B & 91 (73--98) & 35 & 35 & 0.81 (0.55, 1.00) & 1 / 1\\
4 Middleware + version & A & 100 (86--100) & 0 & 0 & undefined & 0 / 0\\
 & B & 83 (63--93) & 0 & 17 & 0.00 (-0.89, 0.89) & 4 / 0\\
5 Algorithms cited & A & 83 (63--93) & 65 & 83 & 0.57 (0.18, 0.95) & 4 / 0\\
 & B & 83 (63--93) & 65 & 74 & 0.59 (0.23, 0.96) & 3 / 1\\
6 Parameters & A & 83 (63--93) & 74 & 91 & 0.43 (-0.09, 0.94) & 4 / 0\\
 & B & 83 (63--93) & 74 & 91 & 0.43 (-0.09, 0.94) & 4 / 0\\
7 Test environment & A & 83 (63--93) & 65 & 57 & 0.64 (0.32, 0.96) & 1 / 3\\
 & B & 87 (68--95) & 65 & 70 & 0.70 (0.39, 1.00) & 2 / 1\\
8 Success metric & A & 91 (73--98) & 74 & 83 & 0.75 (0.41, 1.00) & 2 / 0\\
 & B & 83 (63--93) & 74 & 91 & 0.43 (-0.09, 0.94) & 4 / 0\\
9 Own code/data & A & 100 (86--100) & 13 & 13 & 1.00 (1.00, 1.00) & 0 / 0\\
 & B & 100 (86--100) & 13 & 13 & 1.00 (1.00, 1.00) & 0 / 0\\
SLAM / mapping & A & 83 (63--93) & 48 & 39 & 0.65 (0.34, 0.96) & 1 / 3\\
 & B & 83 (63--93) & 48 & 48 & 0.65 (0.34, 0.96) & 2 / 2\\
Autonomous navigation & A & 74 (54--87) & 43 & 52 & 0.48 (0.12, 0.84) & 4 / 2\\
 & B & 70 (49--84) & 43 & 65 & 0.41 (0.05, 0.78) & 6 / 1\\
Obstacle avoidance & A & 74 (54--87) & 26 & 43 & 0.44 (0.06, 0.83) & 5 / 1\\
 & B & 83 (63--93) & 26 & 35 & 0.59 (0.23, 0.96) & 3 / 1\\
Object detection & A & 87 (68--95) & 30 & 35 & 0.70 (0.39, 1.00) & 2 / 1\\
 & B & 74 (54--87) & 30 & 48 & 0.47 (0.10, 0.83) & 5 / 1\\
Instrument reading & A & 91 (73--98) & 9 & 17 & 0.62 (0.12, 1.00) & 2 / 0\\
 & B & 96 (79--99) & 9 & 13 & 0.78 (0.35, 1.00) & 1 / 0\\
Target tracking & A & 100 (86--100) & 4 & 4 & 1.00 (1.00, 1.00) & 0 / 0\\
 & B & 100 (86--100) & 4 & 4 & 1.00 (1.00, 1.00) & 0 / 0\\
Real-world validation & A & 65 (45--81) & 65 & 83 & 0.13 (-0.35, 0.62) & 6 / 2\\
 & B & 70 (49--84) & 65 & 78 & 0.26 (-0.19, 0.72) & 5 / 2\\
4 Middleware + version (blind codes) & A & 78 (58--90) & 22 & 0 & 0.00 (-0.78, 0.78) & 0 / 5\\
 & B & 78 (58--90) & 22 & 17 & 0.31 (-0.22, 0.85) & 2 / 3\\
2 Sensor model (blind codes) & A & 91 (73--98) & 43 & 35 & 0.82 (0.58, 1.00) & 0 / 2\\
 & B & 91 (73--98) & 43 & 43 & 0.82 (0.59, 1.00) & 1 / 1\\
\bottomrule
\end{tabularx}

}

\section*{Table S5. Class-level validation of the controlled vocabulary}

Agreement between the classes derived from the pipeline's sensor and
compute strings and those derived from the blind human strings, per class,
on the \NcodingStudies{} validation studies. A dash means $\kappa$ is
undefined because neither side assigned the class.

\begin{tabularx}{\textwidth}{llCCCC}
\toprule
\textbf{Field} & \textbf{Class} & \textbf{Agr., \%} & \textbf{$\kappa$} & \textbf{Human positives} & \textbf{Pipeline positives}\\
\midrule
Sensor & LiDAR & 96 & 0.90 & 7 & 8\\
Sensor & Monocular camera & 65 & 0.24 & 4 & 10\\
Sensor & RGB-D / depth & 87 & 0.68 & 6 & 7\\
Sensor & Thermal / IR & 91 & 0.62 & 2 & 4\\
Sensor & Stereo & 100 & 1.00 & 1 & 1\\
Sensor & Event camera & 100 & -- & 0 & 0\\
Compute & Embedded GPU module & 100 & 1.00 & 4 & 4\\
Compute & x86 mini PC & 96 & 0.78 & 2 & 3\\
Compute & Laptop/desktop-class CPU & 96 & 0.65 & 1 & 2\\
Compute & SBC / microcontroller & 96 & 0.65 & 1 & 2\\
Compute & Vendor-integrated computer & 100 & -- & 0 & 0\\
\bottomrule
\end{tabularx}


\section*{Table S6. Sensitivity analysis}

Main results for all readable studies and without the \NfailsCurrent{}
earlier studies that would not meet the eligibility criteria as applied in
round~2. Shares are percentages of the readable studies; differences are in
percentage points.

{\small
\begin{tabularx}{\textwidth}{XCCC}
\toprule
\textbf{Quantity} & \textbf{All ($n=107$)} & \textbf{Without flagged ($n=91$)} & \textbf{Difference}\\
\midrule
Modality: LiDAR & 56.1\,\% & 54.9\,\% & -1.1\\
Modality: Monocular camera & 43.9\,\% & 42.9\,\% & -1.1\\
Modality: RGB-D / depth & 37.4\,\% & 34.1\,\% & -3.3\\
Modality: Thermal / IR & 21.5\,\% & 22.0\,\% & +0.5\\
Modality: Stereo & 8.4\,\% & 6.6\,\% & -1.8\\
Modality: Event camera & 0.9\,\% & 1.1\,\% & +0.2\\
Capability: SLAM / mapping & 54.2\,\% & 53.8\,\% & -0.4\\
Capability: Autonomous navigation & 71.0\,\% & 67.0\,\% & -4.0\\
Capability: Obstacle avoidance & 49.5\,\% & 47.3\,\% & -2.3\\
Capability: Object detection & 44.9\,\% & 46.2\,\% & +1.3\\
Capability: Instrument reading & 8.4\,\% & 8.8\,\% & +0.4\\
Capability: Target tracking & 6.5\,\% & 5.5\,\% & -1.0\\
Capability: real-world validation & 76.6\,\% & 79.1\,\% & +2.5\\
RCI item 1 Platform model & 76.6\,\% & 79.1\,\% & +2.5\\
RCI item 2 Sensor model & 49.5\,\% & 51.6\,\% & +2.1\\
RCI item 3 Compute model & 41.1\,\% & 45.1\,\% & +3.9\\
RCI item 4 Middleware + version & 14.0\,\% & 14.3\,\% & +0.3\\
RCI item 5 Algorithms cited & 75.7\,\% & 75.8\,\% & +0.1\\
RCI item 6 Parameters & 89.7\,\% & 90.1\,\% & +0.4\\
RCI item 7 Test environment & 72.9\,\% & 74.7\,\% & +1.8\\
RCI item 8 Success metric & 76.6\,\% & 84.6\,\% & +8.0\\
RCI item 9 Own code/data & 10.3\,\% & 9.9\,\% & -0.4\\
RCI median & 5.0 & 5.0 & +0.0\\
RCI mean & 5.1 & 5.3 & +0.2\\
Sensor named at model level & 49.5\,\% & 51.6\,\% & +2.1\\
Onboard computer not reported & 53.3\,\% & 48.4\,\% & -4.9\\
\bottomrule
\end{tabularx}

}

%-----------------------------------------------------------------
\renewcommand{\refname}{References for Table S2}
\begin{thebibliography}{999}
\bibitem[Wisth et al.(2019)]{inc:paper014}
Wisth, D.; Camurri, M.; Fallon, M. Robust Legged Robot State Estimation Using Factor Graph Optimization. \emph{IEEE Robotics and Automation Letters} \textbf{2019}, \emph{4}, 4507--4514. \url{https://doi.org/10.1109/lra.2019.2933768}

\bibitem[Ramezani et al.(2020)]{inc:paper019}
Ramezani, M.; Tinchev, G.; Iuganov, E.; Fallon, M. Online LiDAR-SLAM for Legged Robots with Robust Registration and Deep-Learned Loop Closure. In \emph{2020 IEEE International Conference on Robotics and Automation (ICRA)}; 2020; pp. 4158--4164. \url{https://doi.org/10.1109/icra40945.2020.9196769}

\bibitem[Wang et al.(2020)]{inc:paper040}
Wang, Y.; Ramezani, M.; Fallon, M. Actively Mapping Industrial Structures with Information Gain-Based Planning on a Quadruped Robot. In \emph{2020 IEEE International Conference on Robotics and Automation (ICRA)}; 2020; pp. 8609--8615. \url{https://doi.org/10.1109/icra40945.2020.9197153}

\bibitem[Miller et al.(2020)]{inc:paper005}
Miller, I.D.; Cladera, F.; Cowley, A.; Shivakumar, S.S.; Lee, E.S.; Jarin-Lipschitz, L.; Bhat, A.; Rodrigues, N.; Zhou, A.; Cohen, A.; et al. Mine Tunnel Exploration Using Multiple Quadrupedal Robots. \emph{IEEE Robotics and Automation Letters} \textbf{2020}, \emph{5}, 2840--2847. \url{https://doi.org/10.1109/lra.2020.2972872}

\bibitem[Xin et al.(2020)]{inc:paper007}
Xin, G.; Tiseo, C.; Wolfslag, W.; Smith, J.; Cebe, O.; Li, Z.; Vijayakumar, S.; Mistry, M. Variable Autonomy of Whole-body Control for Inspection and Intervention in Industrial Environments using Legged Robots. In \emph{2020 IEEE 16th International Conference on Automation Science and Engineering (CASE)}; 2020; pp. 1415--1420. \url{https://doi.org/10.1109/case48305.2020.9216813}

\bibitem[Brandao et al.(2020)]{inc:paper073}
Brandao, M.; Aladag, O.B.; Havoutis, I. GaitMesh: Controller-Aware Navigation Meshes for Long-Range Legged Locomotion Planning in Multi-Layered Environments. \emph{IEEE Robotics and Automation Letters} \textbf{2020}, \emph{5}, 3596--3603. \url{https://doi.org/10.1109/lra.2020.2979628}

\bibitem[Buchanan et al.(2020)]{inc:r2-0477}
Buchanan, R.; Camurri, M.; Fallon, M. Haptic Sequential Monte Carlo Localization for Quadrupedal Locomotion in Vision-Denied Scenarios. In \emph{2020 IEEE/RSJ International Conference on Intelligent Robots and Systems (IROS)}; 2020; pp. 3657--3663. \url{https://doi.org/10.1109/iros45743.2020.9341128}

\bibitem[Bouman et al.(2020)]{inc:r2-0192}
Bouman, A.; Ginting, M.F.; Alatur, N.; Palieri, M.; Fan, D.D.; Touma, T.; Pailevanian, T.; Kim, S.K.; Otsu, K.; Burdick, J.; et al. Autonomous Spot: Long-Range Autonomous Exploration of Extreme Environments with Legged Locomotion. In \emph{2020 IEEE/RSJ International Conference on Intelligent Robots and Systems (IROS)}; 2020; pp. 2518--2525. \url{https://doi.org/10.1109/iros45743.2020.9341361}

\bibitem[Kim et al.(2021)]{inc:paper065}
Kim, J.; Chung, D.; Kim, Y.; Kim, H. Semantic segmentation of 3D point cloud data acquired from robot dog for scaffold monitoring. In \emph{Proceedings of the International Symposium on Automation and Robotics in Construction (IAARC)}; 2021. \url{https://doi.org/10.22260/isarc2021/0106}

\bibitem[Jud et al.(2021)]{inc:paper087}
Jud, D.; Kerscher, S.; Wermelinger, M.; Jelavic, E.; Egli, P.; Leemann, P.; Hottiger, G.; Hutter, M. HEAP - The autonomous walking excavator. \emph{Automation in Construction} \textbf{2021}, \emph{129}, 103783. \url{https://doi.org/10.1016/j.autcon.2021.103783}

\bibitem[Dong et al.(2021)]{inc:paper075}
Dong, X.; Zhang, B.; Hu, X.; Meng, J.; Dong, J. Design of Quadruped Inspection Robot for Substation. In \emph{Journal of Physics: Conference Series}; 2021; pp. 052010. \url{https://doi.org/10.1088/1742-6596/1885/5/052010}

\bibitem[Wu et al.(2022)]{inc:paper036}
Wu, C.; Zhou, Y.; Zhang, Y.; Li, H.; Wang, X.; Li, Z.; Xu, Y.; Ni, P. Analysis on the locomotion of cable tunnel inspection quadruped robot based on deep reinforcement learning. In \emph{IET Conference Proceedings}; 2022; pp. 1995--1999. \url{https://doi.org/10.1049/icp.2022.0457}

\bibitem[Mattamala et al.(2021)]{inc:paper038}
Mattamala, M.; Ramezani, M.; Camurri, M.; Fallon, M. Learning Camera Performance Models for Active Multi-Camera Visual Teach and Repeat. In \emph{2021 IEEE International Conference on Robotics and Automation (ICRA)}; 2021; pp. 14346--14352. \url{https://doi.org/10.1109/icra48506.2021.9561650}

\bibitem[Stachowiak et al.(2021)]{inc:paper078}
Stachowiak, M.; Koperska, W.; Stefaniak, P.; Skoczylas, A.; Anufriiev, S. Procedures of Detecting Damage to a Conveyor Belt with Use of an Inspection Legged Robot for Deep Mine Infrastructure. \emph{Minerals} \textbf{2021}, \emph{11}, 1040. \url{https://doi.org/10.3390/min11101040}

\bibitem[Skoczylas et al.(2021)]{inc:paper085}
Skoczylas, A.; Stefaniak, P.; Anufriiev, S.; Jachnik, B. Belt Conveyors Rollers Diagnostics Based on Acoustic Signal Collected Using Autonomous Legged Inspection Robot. \emph{Applied Sciences} \textbf{2021}, \emph{11}, 2299. \url{https://doi.org/10.3390/app11052299}

\bibitem[Gehring et al.(2021)]{inc:r2-1437}
Gehring, C.; Fankhauser, P.; Isler, L.; Diethelm, R.; Bachmann, S.; Potz, M.; Gerstenberg, L.; Hutter, M. ANYmal in the Field: Solving Industrial Inspection of an Offshore HVDC Platform with a Quadrupedal Robot. In \emph{Springer Proceedings in Advanced Robotics}; 2021; pp. 247--260. \url{https://doi.org/10.1007/978-981-15-9460-1_18}

\bibitem[Halder et al.(2022)]{inc:paper069}
Halder, S.; Afsari, K.; Serdakowski, J.; DeVito, S.; Ensafi, M.; Thabet, W. Real-Time and Remote Construction Progress Monitoring with a Quadruped Robot Using Augmented Reality. \emph{Buildings} \textbf{2022}, \emph{12}, 2027. \url{https://doi.org/10.3390/buildings12112027}

\bibitem[Wetzel et al.(2022)]{inc:paper070}
Wetzel, E.M.; Liu, J.; Leathem, T.; Sattineni, A. The Use of Boston Dynamics SPOT in Support of LiDAR Scanning on Active Construction Sites. In \emph{Proceedings of the International Symposium on Automation and Robotics in Construction (IAARC)}; 2022. \url{https://doi.org/10.22260/isarc2022/0014}

\bibitem[Ding et al.(2022)]{inc:paper018}
Ding, W.; Zhu, X.; Guo, T.; Liu, X.; Zhang, Y. Research on Quadruped Robot Inspection System Based on Digital Twin Technology. In \emph{2022 2nd International Conference on Electrical Engineering and Control Science (IC2ECS)}; 2022; pp. 1256--1260. \url{https://doi.org/10.1109/ic2ecs57645.2022.10088013}

\bibitem[Kulkarni et al.(2022)]{inc:paper013}
Kulkarni, M.; Dharmadhikari, M.; Tranzatto, M.; Zimmermann, S.; Reijgwart, V.; De Petris, P.; Nguyen, H.; Khedekar, N.; Papachristos, C.; Ott, L.; et al. Autonomous Teamed Exploration of Subterranean Environments using Legged and Aerial Robots. In \emph{2022 International Conference on Robotics and Automation (ICRA)}; 2022; pp. 3306--3313. \url{https://doi.org/10.1109/icra46639.2022.9812401}

\bibitem[Mattamala et al.(2022)]{inc:paper025}
Mattamala, M.; Chebrolu, N.; Fallon, M. An Efficient Locally Reactive Controller for Safe Navigation in Visual Teach and Repeat Missions. \emph{IEEE Robotics and Automation Letters} \textbf{2022}, \emph{7}, 2353--2360. \url{https://doi.org/10.1109/lra.2022.3143196}

\bibitem[Tang et al.(2022)]{inc:paper067}
Tang, W.; Xue, C.; Li, C.; Zhu, Q. Towards Coordinated Multi-Robot Exploration under Bandwidth-constrained Conditions. In \emph{2022 IEEE International Symposium on Safety, Security, and Rescue Robotics (SSRR)}; 2022; pp. 180--187. \url{https://doi.org/10.1109/ssrr56537.2022.10018691}

\bibitem[Lysakowski et al.(2022)]{inc:r2-0387}
Lysakowski, M.; Nowicki, M.R.; Buchanan, R.; Camurri, M.; Fallon, M.; Walas, K. Unsupervised Learning of Terrain Representations for Haptic Monte Carlo Localization. In \emph{2022 International Conference on Robotics and Automation (ICRA)}; 2022; pp. 4642--4648. \url{https://doi.org/10.1109/icra46639.2022.9812296}

\bibitem[Jang et al.(2021)]{inc:r2-1156}
Jang, Y.; Seol, W.; Lee, K.; Kim, K.; Kim, S. Development of quadruped robot for inspection of underground pipelines in nuclear power plants. \emph{Electronics Letters} \textbf{2021}, \emph{58}, 234--236. \url{https://doi.org/10.1049/ell2.12414}

\bibitem[He et al.(2023)]{inc:paper026}
He, Y.; Cheng, L.; Pan, Y.; Li, Y.; Wang, D.; Zheng, H. Multi-sensory Olfactory Quadruped Robot for Odor Source Localization*. In \emph{2022 IEEE International Conference on Cyborg and Bionic Systems (CBS)}; 2023; pp. 332--335. \url{https://doi.org/10.1109/cbs55922.2023.10115389}

\bibitem[Molnar et al.(2023)]{inc:paper084}
Molnar, T.G.; Tighe, K.; Ubellacker, W.; Kalantari, A.; Ames, A.D. Mechanical Design, Planning, and Control for Legged Robots in Distillation Columns. \emph{Journal of Computational and Nonlinear Dynamics} \textbf{2023}, \emph{18}, 061001. \url{https://doi.org/10.1115/1.4056659}

\bibitem[Halder et al.(2023)]{inc:paper066}
Halder, S.; Afsari, K.; Chiou, E.; Patrick, R.; Hamed, K.A. Construction inspection \&amp; monitoring with quadruped robots in future human-robot teaming: A preliminary study. \emph{Journal of Building Engineering} \textbf{2023}, \emph{65}, 105814. \url{https://doi.org/10.1016/j.jobe.2022.105814}

\bibitem[Fiorucci et al.(2023)]{inc:paper080}
Fiorucci, M.; Schillaci, G.; Tannous, M.; Bianchi, L.; Salusti, L.; Cioncolini, S.; Politano, A.; De Magistris, G. Evaluating the Robustness of Autonomous Inspections in the Energy Industry with a Quadruped Robot. In \emph{ADIPEC}; 2023; pp. D031S118R005. \url{https://doi.org/10.2118/217088-ms}

\bibitem[Li et al.(2022)]{inc:paper077}
Li, L.; Wei, Z.; Zhang, T.; Cai, D. Three‐dimensional dead reckoning of wall‐climbing robot based on information fusion of compound extended Kalman filter. \emph{Journal of Field Robotics} \textbf{2022}, \emph{40}, 505--520. \url{https://doi.org/10.1002/rob.22144}

\bibitem[Zhang et al.(2023)]{inc:paper027}
Zhang, Q.; Cheng, J.; Hu, X.; Yu, Y.; Wu, Q. Enhancing Industrial Safety: A Comprehensive Approach to Detecting and Mitigating Potential Hazards. In \emph{2023 China Automation Congress (CAC)}; 2023; pp. 8279--8284. \url{https://doi.org/10.1109/cac59555.2023.10451937}

\bibitem[Ly et al.(2023)]{inc:paper044}
Ly, K.T.; Munks, M.; Merkt, W.; Havoutis, I. Asymptotically Optimized Multi-Surface Coverage Path Planning for Loco-Manipulation in Inspection and Monitoring. In \emph{2023 IEEE 19th International Conference on Automation Science and Engineering (CASE)}; 2023; pp. 1--7. \url{https://doi.org/10.1109/case56687.2023.10260625}

\bibitem[Crespo et al.(2023)]{inc:paper064}
Crespo, C.; Rodríguez, F. Integration of robotics in underground mining construction works. \emph{Expanding Underground - Knowledge and Passion to Make a Positive Impact on the World} \textbf{2023}, 2414--2421. \url{https://doi.org/10.1201/9781003348030-290}

\bibitem[Chu et al.(2023)]{inc:paper072}
Chu, H.; Cao, R.; Deng, L. A collaborative inspection system composed of quadruped and flying robot for crack segmentation in tunnel environment. \emph{Expanding Underground - Knowledge and Passion to Make a Positive Impact on the World} \textbf{2023}, 2406--2413. \url{https://doi.org/10.1201/9781003348030-289}

\bibitem[Brogaard et al.(2023)]{inc:r2-0495}
Brogaard, R.Y.; Hewitt, R.A.; Etter, S.; Kalantari, A.; Boukas, E. Absolute Localization in Feature-Poor Industrial Confined Spaces. In \emph{2023 IEEE International Symposium on Safety, Security, and Rescue Robotics (SSRR)}; 2023; pp. 69--75. \url{https://doi.org/10.1109/ssrr59696.2023.10499935}

\bibitem[Lee et al.(2024)]{inc:paper021}
Lee, J.; Kim, J.; Ames, A.D. Safety-critical Autonomous Inspection of Distillation Columns using Quadrupedal Robots Equipped with Roller Arms. In \emph{2024 IEEE/RSJ International Conference on Intelligent Robots and Systems (IROS)}; 2024; pp. 9094--9101. \url{https://doi.org/10.1109/iros58592.2024.10802679}

\bibitem[Lee et al.(2024)]{inc:paper068}
Lee, J.; Kim, J.; Ubellacker, W.; Molnar, T.G.; Ames, A.D. Safety-critical Control of Quadrupedal Robots with Rolling Arms for Autonomous Inspection of Complex Environments. In \emph{2024 IEEE International Conference on Robotics and Automation (ICRA)}; 2024; pp. 3485--3491. \url{https://doi.org/10.1109/icra57147.2024.10610504}

\bibitem[Dassori et al.(2024)]{inc:paper006}
Dassori, I.; Adams, M.; Vasquez, J. Four-Legged Gait Control via the Fusion of Computer Vision and Reinforcement Learning. In \emph{2024 27th International Conference on Information Fusion (FUSION)}; 2024; pp. 1--8. \url{https://doi.org/10.23919/fusion59988.2024.10706406}

\bibitem[Gan et al.(2024)]{inc:paper086}
Gan, V.J.; Hu, D.; Zhai, R.; Wang, Y. Automated As-built 3D Reconstruction Using Quadruped Robot and 3D LiDAR Sensor. In \emph{Proceedings of the International Symposium on Automation and Robotics in Construction (IAARC)}; 2024. \url{https://doi.org/10.22260/isarc2024/0018}

\bibitem[Baru A(2024)]{inc:paper079}
Baru A Exploring the use of Quadruped Robots in Construction Environments: The Case of Small Item Transportation. In \emph{International Conference On Construction In The 21St Century}; 2024.

\bibitem[Wu et al.(2024)]{inc:paper063}
Wu, Y.; Wang, Q.; Xiao, Y.; Song, Y.; Mao, W.; Wang, P. Path planning of quadruped robot for urban natural gas pipe leakage inspection based on optimized RRT* and DWA algorithms. \emph{International Journal of Metrology and Quality Engineering} \textbf{2024}, \emph{15}, 18. \url{https://doi.org/10.1051/ijmqe/2024017}

\bibitem[Tsenis et al.(2024)]{inc:paper081}
Tsenis, T.; Tromaras, A.; Kappatos, V. NDE Using Quadruped Robotic Platform for Renewable Energy Producing Infrastructures. In \emph{2024 IEEE 8th Energy Conference (ENERGYCON)}; 2024; pp. 1--6. \url{https://doi.org/10.1109/energycon58629.2024.10488808}

\bibitem[Aziz et al.(2024)]{inc:paper062}
Aziz, A.; Herbers, P.; Bayer, H.; König, M. Fully Autonomous Fire Safety Equipment Inspection Missions on a Legged Robot. In \emph{Computing in Civil Engineering 2023}; 2024; pp. 804--812. \url{https://doi.org/10.1061/9780784485224.097}

\bibitem[Ginting et al.(2024)]{inc:paper071}
Ginting, M.F.; Fan, D.D.; Kim, S.K.; Kochenderfer, M.J.; Agha-Mohammadi, A.A. Semantic Belief Behavior Graph: Enabling Autonomous Robot Inspection in Unknown Environments. In \emph{2024 IEEE/RSJ International Conference on Intelligent Robots and Systems (IROS)}; 2024; pp. 7604--7610. \url{https://doi.org/10.1109/iros58592.2024.10802336}

\bibitem[Hu et al.(2025)]{inc:paper083}
Hu, X.; Assaad, R.H. Developing a Quadruped Robot-Based 3D and Panoramic Thermal Mapping System: Moving towards Smart and Automated Building Thermal Inspection Using LiDAR-Based Digital Twins and Sensor Fusion Technologies. In \emph{Computing in Civil Engineering 2024}; 2025; pp. 774--783. \url{https://doi.org/10.1061/9780784486115.082}

\bibitem[Rowell et al.(2024)]{inc:paper041}
Rowell, J.; Zhang, L.; Fallon, M. LiSTA: Geometric Object-Based Change Detection in Cluttered Environments. In \emph{2024 IEEE International Conference on Robotics and Automation (ICRA)}; 2024; pp. 3632--3638. \url{https://doi.org/10.1109/icra57147.2024.10610102}

\bibitem[Chang et al.(2024)]{inc:paper012}
Chang, J.; Park, J. Design of a Robot System for Surveillance and Anomaly Detection in Industrial Environments. In \emph{2024 15th International Conference on Information and Communication Technology Convergence (ICTC)}; 2024; pp. 2229--2234. \url{https://doi.org/10.1109/ictc62082.2024.10827327}

\bibitem[Zhou et al.(2024)]{inc:paper015}
Zhou, Z.; Huang, P.; Zheng, S.; Xu, Z.; Zhuang, Z. Proximity Estimation with Position Adjustment for Autonomous Industrial Inspection. In \emph{2024 International Conference on Artificial Intelligence of Things and Systems (AIoTSys)}; 2024; pp. 1--6. \url{https://doi.org/10.1109/aiotsys63104.2024.10780707}

\bibitem[Reese et al.(2024)]{inc:paper037}
Reese, R.; Kovarovics, A.; Charles, A.; Koduru, C.; Tanveer, M.H.; Voicu, R.C. Optimizing Data Capture Through Object Recognition for Efficient Sensor and Camera Management with a Quadruped Robot. In \emph{SoutheastCon 2024}; 2024; pp. 1125--1130. \url{https://doi.org/10.1109/southeastcon52093.2024.10500066}

\bibitem[Zagorski et al.(2024)]{inc:paper082}
Zagorski, M.; Miltchev, R. FACILITATION OF THE 3D SCANNING PROCESS OF INDUSTRIAL SITES USING A SELF-MOVING AUTONOMOUS ROBOTIC SYSTEM. In \emph{ENVIRONMENT. TECHNOLOGY. RESOURCES. Proceedings of the International Scientific and Practical Conference}; 2024; pp. 348--351. \url{https://doi.org/10.17770/etr2024vol3.8131}

\bibitem[Hou et al.(2024)]{inc:paper039}
Hou, E.; Zhao, R. VertiPaw: Development of a Vertical Climbing Robot with Adaptive Suction and Suspension Foot Design for Infastructure Monitoring. In \emph{2024 Eighth IEEE International Conference on Robotic Computing (IRC)}; 2024; pp. 116--123. \url{https://doi.org/10.1109/irc63610.2024.00027}

\bibitem[Turrisi et al.(2024)]{inc:paper074}
Turrisi, G.; Schulze, L.; Medeiros, V.S.; Semini, C.; Barasuol, V. PACC: A Passive-Arm Approach for High-Payload Collaborative Carrying with Quadruped Robots Using Model Predictive Control. In \emph{2024 IEEE/RSJ International Conference on Intelligent Robots and Systems (IROS)}; 2024; pp. 11139--11146. \url{https://doi.org/10.1109/iros58592.2024.10801456}

\bibitem[Wei et al.(2024)]{inc:paper033}
Wei, Z.; Liu, S.; Song, R. A Path Planning Method for Quadruped Robot Based on 3D Elevation Diagram. In \emph{2024 7th International Symposium on Autonomous Systems (ISAS)}; 2024; pp. 1--6. \url{https://doi.org/10.1109/isas61044.2024.10552499}

\bibitem[Sun et al.(2024)]{inc:paper076}
Sun, W.; Shen, D.; Liu, K.; Li, Z.; Zhou, B. Multi-sensor Fusion-based Perception and Navigation Framework for Quadruped Robots in Complex Environments*. In \emph{2024 43rd Chinese Control Conference (CCC)}; 2024; pp. 4561--4566. \url{https://doi.org/10.23919/ccc63176.2024.10662127}

\bibitem[Chen et al.(2024)]{inc:paper048}
Chen, L.; Peng, P.; Wang, L. MAMRS: Mining Automatic Meter Reading System Based on Improved Deep Learning Algorithm Using Quadruped Robots. \emph{Applied Sciences} \textbf{2024}, \emph{14}, 10949. \url{https://doi.org/10.3390/app142310949}

\bibitem[Balan et al.(2024)]{inc:paper035}
Balan, K.; Li, P.Y.; Bandstra, M.S.; Patel, K.; Rofors, E.; Salathe, M.; MacDonald, T.D.; Park, K.; Godfrey, E.; Quiter, B.J. Robot Path Planning Utilizing Object Recognition for Inspection of Nuclear Material Containers. In \emph{2024 IEEE International Symposium on Safety Security Rescue Robotics (SSRR)}; 2024; pp. 46--51. \url{https://doi.org/10.1109/ssrr62954.2024.10770038}

\bibitem[Xu et al.(2024)]{inc:paper049}
Xu, B.; Zhang, X.; Yu, X.; Ou, Y.; Zhang, K.; Cai, H.; Zhao, J.; Fan, J. A Motion Planner Based on Mask-D3QN of Quadruped Robot Motion for Steam Generator. \emph{Biomimetics} \textbf{2024}, \emph{9}, 592. \url{https://doi.org/10.3390/biomimetics9100592}

\bibitem[Cheng et al.(2024)]{inc:paper023}
Cheng, J.; Zhang, M.; Ma, L.; Chen, H.; Gan, Y.; Huang, D. A Hybrid-Dimensional Laser SLAM Framework for Indoor Quadruped Inspection Robots. \emph{IEEE Sensors Journal} \textbf{2024}, \emph{24}, 16935--16942. \url{https://doi.org/10.1109/jsen.2024.3382336}

\bibitem[Li et al.(2024)]{inc:r2-0771}
Li, J.; Liu, Z.; Li, S.; Jiang, J.; Li, Y.; Tian, C.; Wang, G. Motion planning for a quadruped robot in heat transfer tube inspection. \emph{Automation in Construction} \textbf{2024}, \emph{168}, 105753. \url{https://doi.org/10.1016/j.autcon.2024.105753}

\bibitem[Baird et al.(2024)]{inc:r2-1052}
Baird, C.; Nokleby, S. Autonomous firefighting using a quadruped robot. \emph{Transactions of the Canadian Society for Mechanical Engineering} \textbf{2024}, \emph{48}, 605--616. \url{https://doi.org/10.1139/tcsme-2023-0175}

\bibitem[Huang et al.(2024)]{inc:r2-1239}
Huang, Y.; Chen, Z.; Chen, J.; Wen, D.; Wang, Y. Optimal path planning for large-scale emergency inspection of quadruped robot considering complex terrain. In \emph{Proceedings of the 2024 3rd International Symposium on Control Engineering and Robotics}; 2024; pp. 414--419. \url{https://doi.org/10.1145/3679409.3679486}

\bibitem[Moon et al.(2024)]{inc:r2-1169}
Moon, S.; Kim, S.K.; Peltzer, O.; Kochenderfer, M.J.; Khattak, S. Efficient Line-of-Sight Viewpoint Sampling in Complex Environments for Autonomous Surface Inspection. In \emph{2024 IEEE International Symposium on Safety Security Rescue Robotics (SSRR)}; 2024; pp. 198--203. \url{https://doi.org/10.1109/ssrr62954.2024.10770030}

\bibitem[Damigos et al.(2026)]{inc:paper042}
Damigos, G.; Pagliari, E.; Sandberg, S.; Nikolakopoulos, G. 5G-Enabled Robots: Differentiated Connectivity for Varying Mission Requirements Through Dynamic QoS. \emph{IEEE Transactions on Vehicular Technology} \textbf{2026}, \emph{75}, 9352--9367. \url{https://doi.org/10.1109/tvt.2025.3646562}

\bibitem[Gan et al.(2025)]{inc:paper047}
Gan, V.J.; Hu, D.; Wang, Y.; Zhai, R. Automated indoor 3D scene reconstruction with decoupled mapping using quadruped robot and LiDAR sensor. \emph{Computer-Aided Civil and Infrastructure Engineering} \textbf{2025}, \emph{40}, 2209--2226. \url{https://doi.org/10.1111/mice.13450}

\bibitem[Zeng et al.(2025)]{inc:paper003}
Zeng, Y.; Song, Z.; Tian, S.; Cheng, G. Integrated Quadruped–UAV Robotic Platform for Efficient Construction Site Inspection. In \emph{2025 5th International Conference on Conference on Robotics, Automation and Intelligent Control (ICRAIC)}; 2025; pp. 1--5. \url{https://doi.org/10.1109/icraic67376.2025.11376350}

\bibitem[Scheidemann et al.(2025)]{inc:paper030}
Scheidemann, C.; Werner, L.; Reijgwart, V.; Cramariuc, A.; Chomarat, J.; Chiu, J.R.; Siegwart, R.; Hutter, M. Obstacle-Avoidant Leader Following with a Quadruped Robot. In \emph{2025 IEEE International Conference on Robotics and Automation (ICRA)}; 2025; pp. 1407--1413. \url{https://doi.org/10.1109/icra55743.2025.11127912}

\bibitem[Chung et al.(2025)]{inc:paper053}
Chung, D.; Kim, J.; Paik, S.; Im, S.; Kim, H. Automated system of scaffold point cloud data acquisition using a robot dog. \emph{Automation in Construction} \textbf{2025}, \emph{170}, 105944. \url{https://doi.org/10.1016/j.autcon.2024.105944}

\bibitem[Stührenberg et al.(2025)]{inc:paper057}
Stührenberg, J.; Smarsly, K. LIO-BIM – Coupling lidar inertial odometry with building information modeling for robot localization and mapping. \emph{Advanced Engineering Informatics} \textbf{2025}, \emph{66}, 103477. \url{https://doi.org/10.1016/j.aei.2025.103477}

\bibitem[Kim et al.(2025)]{inc:paper061}
Kim, J.; Chung, D.; Kim, J.; Kim, H. Integrated Framework for Scaffold Monitoring: From Point Cloud Acquisition to Anomaly Detection. In \emph{Proceedings of the International Symposium on Automation and Robotics in Construction (IAARC)}; 2025. \url{https://doi.org/10.22260/isarc2025/0192}

\bibitem[Xu et al.(2025)]{inc:paper011}
Xu, B.; Lang, Y.; Chen, Q.; Yin, L.; Li, Q. A Motion Planning Method for Tube Plate Quadruped Robot Based on Target-Oriented D* Search. In \emph{2025 IEEE International Conference on Robotics and Biomimetics (ROBIO)}; 2025; pp. 1152--1157. \url{https://doi.org/10.1109/robio66223.2025.11377431}

\bibitem[Kim et al.(2025)]{inc:paper058}
Kim, H.B.; Kim, C.; Ham, B.I.; Kang, J.; Choi, M.; Choi, K.H.; Kim, K.S. Development of Remote Piping Inspection System with Dual-Mode Locomotion Quadruped Robot. \emph{Lecture Notes in Networks and Systems} \textbf{2025}, 307--319. \url{https://doi.org/10.1007/978-3-031-92011-0_25}

\bibitem[Xing et al.(2025)]{inc:paper028}
Xing, F.; Sun, C.; Feng, W.; Zhang, J.; Liu, H.; Wang, W.; Zhang, R.; Li, D. Research on Core Technology and Field Application of Electric Power Generation of Quadruped Robot (Robot Dog). In \emph{2025 IEEE PES Innovative Smart Grid Technologies - Asia (ISGT Asia)}; 2025; pp. 406--411. \url{https://doi.org/10.1109/isgtasia63446.2025.11431373}

\bibitem[Mu et al.(2025)]{inc:paper043}
Mu, Y.; Zhu, H.; Du, Y.; Tang, X. DP-SLAM: A Tightly-Coupled SLAM Method with Degradation Perception. In \emph{2025 China Automation Congress (CAC)}; 2025; pp. 2500--2505. \url{https://doi.org/10.1109/cac67268.2025.11486644}

\bibitem[Hu et al.(2025)]{inc:paper010}
Hu, L.; Ding, L.; Yang, H.; Liu, T.; Zhang, A.; Chen, S.; Gao, H.; Xu, P.; Deng, Z. LNO-Driven Deep RL-MPC: Hierarchical Adaptive Control Architecture for Dynamic Legged Locomotion. \emph{IEEE Transactions on Industrial Informatics} \textbf{2025}, \emph{21}, 8574--8584. \url{https://doi.org/10.1109/tii.2025.3584522}

\bibitem[Yang et al.(2025)]{inc:paper016}
Yang, H.; Tao, R.; Zhang, S.; Chen, L.; Zhang, T.; Liu, M.; Wang, H.; Wu, Y. Closed-Loop Safety Inspection Application via Embodied Robots. In \emph{2025 2nd International Symposium on AI and Cybersecurity (ISAICS)}; 2025; pp. 1--5. \url{https://doi.org/10.1109/isaics66888.2025.11350245}

\bibitem[Zhu et al.(2025)]{inc:paper020}
Zhu, S.; Mou, L.; Li, D.; Ye, B.; Huang, R.; Zhao, H. VR-Robo: A Real-to-Sim-to-Real Framework for Visual Robot Navigation and Locomotion. \emph{IEEE Robotics and Automation Letters} \textbf{2025}, \emph{10}, 7875--7882. \url{https://doi.org/10.1109/lra.2025.3575648}

\bibitem[Mendoza et al.(2025)]{inc:paper024}
Mendoza, E.; Saeteros, J.; Hurel, J.; Medrano, J.; Moon, H.; Quero, L.C.; Yumbla, E.Q.; Yumbla, F. Control System for Center of Mass Position and Angle in Quadruped Robots. In \emph{2025 22nd International Conference on Ubiquitous Robots (UR)}; 2025; pp. 398--404. \url{https://doi.org/10.1109/ur65550.2025.11078136}

\bibitem[Khatiri et al.(2025)]{inc:paper031}
Khatiri, S.; Barrientos, F.E.V.; Wulf, M.; Tonella, P.; Panichella, S. Bridging Research and Practice in Simulation-based Testing of Industrial Robot Navigation Systems. In \emph{2025 40th IEEE/ACM International Conference on Automated Software Engineering (ASE)}; 2025; pp. 3357--3368. \url{https://doi.org/10.1109/ase63991.2025.00277}

\bibitem[Rizvi et al.(2025)]{inc:paper029}
Rizvi, D.; Paul, G.; Roshan, M.C.; Gunatilake, A. Toward Plug-and-Play Asset Management: Spot-Based Visual Surveying in Untagged Industrial Spaces. In \emph{2025 IEEE International Conference on Robotics and Biomimetics (ROBIO)}; 2025; pp. 1378--1383. \url{https://doi.org/10.1109/robio66223.2025.11379746}

\bibitem[DiBattista et al.(2025)]{inc:paper051}
DiBattista, M.A.; Frericks, J.; Garcia, C.I.; Camelio, J.A. SCOUT: an autonomous UHF RFID-equipped robot dog for flexible inventory monitoring. \emph{Manufacturing Letters} \textbf{2025}, \emph{44}, 1525--1532. \url{https://doi.org/10.1016/j.mfglet.2025.06.172}

\bibitem[Fan et al.(2025)]{inc:paper004}
Fan, J.; Xu, J.; Wang, J.; Wang, Y.; Li, X. Application of Multi-Scale Path Planning for Quadruped Robots in Special Environments. In \emph{2025 5th International Conference on Computer, Control and Robotics (ICCCR)}; 2025; pp. 38--42. \url{https://doi.org/10.1109/icccr65461.2025.11072602}

\bibitem[Gao et al.(2025)]{inc:paper017}
Gao, Y.; Awuah-Offei, K. Navigation in Underground Mine Environments: A Simulation Framework for Quadruped Robots. In \emph{2025 IEEE 21st International Conference on Automation Science and Engineering (CASE)}; 2025; pp. 1464--1469. \url{https://doi.org/10.1109/case58245.2025.11163807}

\bibitem[Li et al.(2025)]{inc:paper009}
Li, J.; Liu, Z.; Li, Y.; Li, Y.; Huang, Y.; Wang, G. Localized Coverage Planning for a Heat Transfer Tube Inspection Robot. \emph{IEEE Robotics and Automation Letters} \textbf{2025}, \emph{10}, 3916--3923. \url{https://doi.org/10.1109/lra.2025.3547675}

\bibitem[Staniaszek et al.(2025)]{inc:paper045}
Staniaszek, M.; Brudermüller, L.; You, Y.; Bhattacharyya, R.; Lacerda, B.; Hawes, N. Time-bounded planning with uncertain task duration distributions. \emph{Robotics and Autonomous Systems} \textbf{2025}, \emph{186}, 104926. \url{https://doi.org/10.1016/j.robot.2025.104926}

\bibitem[Liu et al.(2024)]{inc:paper050}
Liu, Z.; Li, J.; Chen, J.; Li, S.; Gu, X.; Jiang, J.; Li, Y.; Yao, S.; Wang, G. The SG‐Climbot: An Adaptable, Efficient, Inspection and Repair Robot for Steam Generator Heat Transfer Tube Sheet. \emph{Journal of Field Robotics} \textbf{2024}, \emph{42}, 1080--1102. \url{https://doi.org/10.1002/rob.22436}

\bibitem[Yao et al.(2025)]{inc:paper001}
Yao, C.; Li, Q.; Yang, S.; Zhao, K.; Guo, B.; Liu, Z.; Jia, L. Cloud-Based Adaptive Remote Control System for Quadruped Robots in Oil and Gas Environments. In \emph{2025 7th International Conference on Intelligent Control, Measurement and Signal Processing (ICMSP)}; 2025; pp. 454--459. \url{https://doi.org/10.1109/icmsp68723.2025.11407652}

\bibitem[Bin Mohamad Nor et al.(2025)]{inc:paper002}
Bin Mohamad Nor, M.H.; Al-Thani Mahmood, I.; Junoh, S.; Omar, M.B.; Bin Ibrahim, R. Real Environment Testing of Quadruped Robot Mobile Manipulator in Oil and Gas Facilities. In \emph{2025 IEEE 6th International Conference in Robotics and Manufacturing Automation (ROMA)}; 2025; pp. 182--187. \url{https://doi.org/10.1109/roma66616.2025.11155452}

\bibitem[Guo et al.(2025)]{inc:paper034}
Guo, J.; Wang, Z.; Bai, W. Learning Quadrupedal Robot Locomotion for Narrow Pipe Inspection. In \emph{2025 IEEE International Conference on Robotics and Biomimetics (ROBIO)}; 2025; pp. 555--561. \url{https://doi.org/10.1109/robio66223.2025.11379856}

\bibitem[Wu et al.(2025)]{inc:paper054}
Wu, Y.; Li, Y.; Xue, W.; Guo, P.; Xiao, T.; Luo, H.; Li, X.; Zhu, X.; Yang, J.; Lin, Z.H. Flexible, wideband triboelectric accelerometer integrated into quadruped robot legs for vibration source detection and localization. \emph{Composites Part B: Engineering} \textbf{2025}, \emph{305}, 112710. \url{https://doi.org/10.1016/j.compositesb.2025.112710}

\bibitem[Zhang et al.(2025)]{inc:paper060}
Zhang, X.; Huang, X.; Hong, K.; Li, Q.; Wang, R.; Zhou, B. TUC-Net: A Point Cloud Segmentation Network Based on Neighborhood Feature Perception Aggregation for Tunnels Under Construction. \emph{IEEE Transactions on Intelligent Transportation Systems} \textbf{2025}, \emph{26}, 10048--10066. \url{https://doi.org/10.1109/tits.2025.3542500}

\bibitem[Shi et al.(2025)]{inc:paper022}
Shi, S.; Mei, M.; Xue, Z.; Li, N.; Miao, Z. Development of a Quadruped Robot Simulator for Warehouse Inspection Based on Unity3D and ROS. In \emph{2025 China Automation Congress (CAC)}; 2025; pp. 5808--5813. \url{https://doi.org/10.1109/cac67268.2025.11486748}

\bibitem[Adams et al.(2025)]{inc:r2-0704}
Adams, J.; Cintas, B.; Sung, N.; Abrahao, A.; Lagos, L.; McDaniel, D. A Behavioral Robotics Approach to Radiation Mapping Using Adaptive Sampling. \emph{Applied Sciences} \textbf{2025}, \emph{15}, 2050. \url{https://doi.org/10.3390/app15042050}

\bibitem[Chen et al.(2025)]{inc:r2-0745}
Chen, Z.; Song, C.; Wang, B.; Tao, X.; Zhang, X.; Lin, F.; Cheng, J.C. Automated reality capture for indoor inspection using BIM and a multi-sensor quadruped robot. \emph{Automation in Construction} \textbf{2025}, \emph{170}, 105930. \url{https://doi.org/10.1016/j.autcon.2024.105930}

\bibitem[Tuomisto et al.(2026)]{inc:paper056}
Tuomisto, J.; Zheng, Y.; Chauhan, I.; Seppänen, O.; Pajarinen, J. Automating on-site object inspection with a quadruped robot and BIM. \emph{Automation in Construction} \textbf{2026}, \emph{181}, 106571. \url{https://doi.org/10.1016/j.autcon.2025.106571}

\bibitem[Lee et al.(2026)]{inc:paper055}
Lee, J.; Kim, Y.; Ko, B.; Kim, T.; Park, J.; Lee, J.; Kim, W. Robust Analog Gauge Reading via Virtual Point-Based Geometric Rectification and P2-YOLO-Pose. \emph{Computer Modeling in Engineering \&amp; Sciences} \textbf{2026}, \emph{147}, 1--10. \url{https://doi.org/10.32604/cmes.2026.080624}

\bibitem[Liu et al.(2026)]{inc:paper088}
Liu, X.; Qu, N.; Zheng, Y.; Wang, H.; He, W. Coordinated motion control of quadruped robot under extreme conditions based on improved MPC. In \emph{IET Conference Proceedings}; 2026; pp. 1081--1089. \url{https://doi.org/10.1049/icp.2025.4629}

\bibitem[Ji et al.(2026)]{inc:paper059}
Ji, X.; Chen, J.; Fu, Y.; Chan, I.; Dong, Z. Text-prompted generative data augmentation and semi-supervised learning for indoor defect detection using quadruped robots. \emph{Expert Systems with Applications} \textbf{2026}, \emph{312}, 131545. \url{https://doi.org/10.1016/j.eswa.2026.131545}

\bibitem[Luo et al.(2026)]{inc:paper032}
Luo, A.; Li, D.; Wang, S.; Kong, S.; Zhu, K.; Yu, J. Perceptive Locomotion and Navigation for Quadruped Robots via Depth-Based Representation. \emph{IEEE Transactions on Automation Science and Engineering} \textbf{2026}, \emph{23}, 8672--8684. \url{https://doi.org/10.1109/tase.2026.3687076}

\bibitem[Prajapati et al.(2026)]{inc:paper008}
Prajapati, D.; Mehra, A.; Kumar, P.; Rana, A.; Surya Prakash, S.K.; Shukla, A. MPPF: Mobile Probe Potential Field for Real-Time Multi-Agent Navigation in Industrial Environments. \emph{IEEE Access} \textbf{2026}, \emph{14}, 11084--11095. \url{https://doi.org/10.1109/access.2026.3655886}

\bibitem[Zhu et al.(2026)]{inc:paper052}
Zhu, C.; Bao, J.; Zhao, Z.; Yang, Z.; Huang, C.; Kang, J.; Xu, H.; Zhang, Z. A digital twin-based quadruped robot system with scene perception, fast communication, and holographic interaction. \emph{Digital Communications and Networks} \textbf{2026}, \emph{12}, 262--272. \url{https://doi.org/10.1016/j.dcan.2025.11.004}

\bibitem[Adams et al.(2026)]{inc:paper046}
Adams, J.; Abrahao, A.; Lagos, L.; McDaniel, D. Autonomous Radiation Mapping Using a Manipulator-Equipped Quadruped with Flexible Behavior Design. \emph{Applied Sciences} \textbf{2026}, \emph{16}, 3500. \url{https://doi.org/10.3390/app16073500}

\bibitem[Aydogmus et al.(2026)]{inc:r2-0348}
Aydogmus, O.; Boztas, G. Zero-Shot Visual Anomaly Detection on a Quadruped Robot Using State of the Art Visual Language Models. \emph{IEEE Access} \textbf{2026}, \emph{14}, 17842--17852. \url{https://doi.org/10.1109/access.2026.3658108}

\bibitem[Ovur et al.(2026)]{inc:r2-0553}
Ovur, E.; Chacón Quesada, R.; Demiris, Y. Optimizing Ergonomics for Robot-to-Human Object Handovers. \emph{IEEE Transactions on Cognitive and Developmental Systems} \textbf{2026}, \emph{18}, 33--42. \url{https://doi.org/10.1109/tcds.2025.3575625}

\bibitem[Son et al.(2026)]{inc:r2-0677}
Son, H.; Tan, K.; Zhang, F. Physics-Guided Robotic Radiation Source Localization Along Arbitrary Measurement Paths in Unstructured Environments. \emph{IEEE Transactions on Automation Science and Engineering} \textbf{2026}, \emph{23}, 14809--14825. \url{https://doi.org/10.1109/tase.2026.3724772}

\bibitem[Feng et al.(2026)]{inc:r2-0740}
Feng, Y.; Feng, S.J.; Zhang, X.L.; Kong, Q.Z.; Zhao, Y. Automatic deformation detection of metro tunnels via point cloud segmentation and geometric analysis. \emph{Automation in Construction} \textbf{2026}, \emph{181}, 106657. \url{https://doi.org/10.1016/j.autcon.2025.106657}

\bibitem[Haage et al.(2026)]{inc:r2-0776}
Haage, M.; Erikshammar, J.; Stehn, L. Toward Frequent Robotic Digitalization of Construction Sites: A Field-Grounded Roadmap. \emph{Buildings} \textbf{2026}, \emph{16}, 2988. \url{https://doi.org/10.3390/buildings16152988}

\bibitem[Qin et al.(2026)]{inc:r2-0934}
Qin, P.; Xu, G.; Wu, Q.; Wei, F.; Zhao, Y.; Zhang, K. SCGait a novel method for person identification applied to legged robots. \emph{Scientific Reports} \textbf{2026}, \emph{16}, 18566. \url{https://doi.org/10.1038/s41598-026-49046-5}

\bibitem[Bayer et al.(2025)]{inc:r2-1004}
Bayer, J.; Faigl, J. Decentralized multi-robot exploration under low-bandwidth communications. \emph{Autonomous Robots} \textbf{2025}, \emph{50}, 7. \url{https://doi.org/10.1007/s10514-025-10234-3}

\bibitem[Chang et al.(2026)]{inc:r2-1031}
Chang, W.Y.; Hung, C.H.; Huang, Z.X.; Chang, H.J. Real-Time Mapping and Planning Intelligent Paths using Optical LiDAR and Quadruped Robot. \emph{The International Archives of the Photogrammetry, Remote Sensing and Spatial Information Sciences} \textbf{2026}, \emph{XLIX-B1-2026}, 383--388. \url{https://doi.org/10.5194/isprs-archives-xlix-b1-2026-383-2026}

\end{thebibliography}

\end{document}
